%==============================================================================
%  Appendix A  Proof of the spectral geometry                      [owner: A1]
%  Source: v8 app:comp1 (v8:701-786).
%  One-statement-one-home: v8's lem:A-coherence and lem:A-gap restated
%  prop:coherence and lem:gap verbatim; both are now proof environments
%  attached to the main-text statements, and the duplicate labels are gone.
%==============================================================================
\section{Proof of the spectral geometry}\label{app:comp1}
% Supports Proposition~\ref{prop:coherence}, Lemma~\ref{lem:gap}, Corollary~\ref{cor:tolerance}.

For the proof of \Cref{lem:gap} we first derive the spectral properties of the
Fiedler vector.

\begin{lemma}[Spectral decomposition of the imbalanced Laplacian]\label{lem:A-decomp}
Under \Cref{asm:clock,asm:cbm}: (i) the smallest eigenvalue is $\lambda_1=0$ with
$v^{(1)}=\tfrac1{\sqrt m}\1_m$; (ii) the Fiedler eigenvalue is $\lambda_2=m\bze$;
(iii) the Fiedler vector is piecewise constant, $\vv_i=+c_A$ for $i\in A$ and
$\vv_i=-c_B$ for $i\in B$, with $c_A=\sqrt{b/(am)}$ and $c_B=\sqrt{a/(bm)}$.
\end{lemma}
\begin{proof}
For any $L=D-S$, $L\1=D\1-S\1=0$, so $\lambda_1=0$ with $v^{(1)}=\tfrac1{\sqrt
m}\1_m$; no assumption on $S$ is needed.

\emph{Step heights.} Assume $\vv$ has the stated form with $c_A,c_B>0$. The
constraints $\norm{\vv}_2=1$ and $\langle v^{(1)},\vv\rangle=0$ give $ac_A=bc_B$,
hence
\[
   a\Bigl(\tfrac{b}{a}c_B\Bigr)^2+bc_B^2=1
   \;\Longrightarrow\;\tfrac{bm}{a}c_B^2=1
   \;\Longrightarrow\; c_B^2=\tfrac{a}{bm},\quad c_A^2=\tfrac{b}{am}.
\]

\emph{Eigenpair verification.} Write $\vv=(c_A\1_a,\,-c_B\1_b)^\top$ with
$ac_A=bc_B$. Using the block form of $L$ and $J_{a\times b}\1_b=b\1_a$,
$J_{b\times a}\1_a=a\1_b$, together with $S_A\1_a=\operatorname{rowsum}(S_A)\1_a$,
\[
   Lv^{(2)}=
   \begin{pmatrix}
     \bigl(D_A-\operatorname{rowsum}(S_A)\bigr)c_A\1_a+b\bze c_B\1_a\\[2pt]
     -a\bze c_A\1_b-\bigl(D_B-\operatorname{rowsum}(S_B)\bigr)c_B\1_b
   \end{pmatrix}.
\]
Since $D_A=\operatorname{rowsum}(S_A)+b\bze$ and
$D_B=\operatorname{rowsum}(S_B)+a\bze$, this reduces to
$\bigl(b\bze(c_A+c_B)\1_a,\,-a\bze(c_A+c_B)\1_b\bigr)^\top$. From $ac_A=bc_B$,
$b(c_A+c_B)=mc_A$ and $a(c_A+c_B)=mc_B$, so $Lv^{(2)}=m\bze\,\vv$ and
$\lambda_2=m\bze$.
\end{proof}

\begin{proof}[Proof of \Cref{prop:coherence}]
Substituting the Fiedler coordinates of \Cref{lem:A-decomp} into
\eqref{eq:coherence_definition},
\[
   \coh=\tfrac m2\max\Bigl\{\tfrac1m+c_A^2,\ \tfrac1m+c_B^2\Bigr\}
       =\tfrac m2\Bigl(\tfrac1m+\tfrac{b}{ma}\Bigr)
       =\tfrac12\Bigl(1+\tfrac ba\Bigr)=\tfrac{1+\eta}{2}. \qedhere
\]
\end{proof}

\begin{proof}[Proof of \Cref{lem:gap}]
The argument uses \Cref{asm:clock,asm:cbm,asm:bounded}. By \Cref{lem:A-decomp} the
Fiedler eigenvalue is $\lambda_2=m\bze=m\Sout^{\max}$.
For $\lambda_3$, standard algebraic-connectivity interlacing gives
$\lambda_3(L)\ge\min\bigl(\lambda_2(L_{C_1}),\lambda_2(L_{C_2})\bigr)$, and applying
the Fiedler-eigenvalue formula to each sub-clan Laplacian gives
$\lambda_2(L_{C_c})\ge n_c\,S_{\mathrm{in}}^{(c)}$. With
$\Sin=\min_c S_{\mathrm{in}}^{(c)}$,
\[
   \min\bigl(\lambda_2(L_{C_1}),\lambda_2(L_{C_2})\bigr)\ge\min(n_1,n_2)\,\Sin=m_{\min}\Sin,
\]
so $\Dlam=\lambda_3-\lambda_2\ge m_{\min}\Sin-m\Sout^{\max}$. Using
$\Sin=\rho+\Sout^{\max}$ and $m=m_{\min}(1+\eta)$,
\[
   \Dlam\ge m_{\min}(\rho+\Sout^{\max})-(1+\eta)m_{\min}\Sout^{\max}
        =m_{\min}\rho-\eta\,m_{\min}\Sout^{\max}. \qedhere
\]
\end{proof}

\begin{proof}[Proof of \Cref{cor:tolerance}]
Positivity of the bound in \Cref{lem:gap} requires
$\eta<\Sin^{\min}/\Sout^{\max}$. The minimum intra-clan similarity is bounded by
the diameter, $\Sin^{\min}\ge e^{-r(\mathcal T)}$, and the maximum cross-clan
similarity by the subtree depth, $\Sout^{\max}\propto e^{-\mathcal O(h(\mathcal T))}$.
Substituting into $\eta<\Sin^{\min}/\Sout^{\max}$ gives \eqref{eq:eta_distance_bound};
the logarithmic scaling $r,h=\mathcal O(\log m)$ under Kingman yields
$\eta\le\mathcal O(e^{\mathcal O(\log m)})=\mathcal O(m^c)$.
\end{proof}
